\section{Discussion}

The empirical findings in the previous section can be read in two ways: as descriptive statistics about a dataset, or as a portrait of a community caught in the process of forming. This section takes the second reading.

\subsection{A community with a visible history}

Most corpora are static snapshots --- the social processes that produced them are invisible in the data. This corpus is unusual in that its own growth structure narrates a social history. The near-zero activity before November 2024 and the sudden ×193 spike in that month are not noise: they mark the clearest structural break in the corpus. A majority of authors first appear in the corpus in November or December 2024, while lower-volume activity is already present from August 2023 onward. The corpus therefore functions simultaneously as a linguistic resource and as a record of a community in transition.

This has an implication for how the corpus should be used. Any NLP model trained on it, or any linguistic analysis conducted with it, is working with language produced under specific social conditions: by people who chose to move platforms, who were actively establishing a new online presence, and who were in many cases reconnecting with each other after a period of platform disruption. Whether these conditions leave traces in the language is an open question, but it is worth asking.

\subsection{Dialogue over broadcasting}

The most striking single finding is the reply rate: 68.3\% of all posts are replies. People are using the platform primarily to talk to each other, not only to publish at an audience.

This is reinforced by the broadcaster/conversationalist split. The small number of institutional accounts (sports news, environmental organisations, community groups) post exclusively original content. Among individual users the reply rate is often above 80--90\%. If institutional accounts are set aside, the community looks even more strongly conversational.

One interpretation is that this reflects the early-adopter character of the community: a small group of people who know each other, or quickly come to know each other, and who use the platform as a space for ongoing dialogue rather than mass communication. As the community grows, this dynamic may shift.

\subsection{Small but coherent}

With 432 authors and 141,013 posts, the Slovene Bluesky community is small by the standards of any major language platform. The Gini coefficient of 0.82 and the fact that 22 authors account for half the corpus confirm that a core group of highly active users sets the tone. The presence of a distinct SUP paddleboarding subgroup --- visible through hashtag patterns --- and the clear daily rhythm consistent with Slovenian working hours both suggest that, despite its small size, the community has genuine internal structure and is not simply a random collection of Slovene speakers who happen to be on the same platform.

\subsection{ATProto as a research infrastructure}

From a methodological standpoint, the experience of building this corpus confirms that ATProto's open design has real consequences for research. The absence of API rate limits or access agreements meant that a full backfill of the discovered Slovene-linked authors was achievable without institutional infrastructure. This stands in contrast to the conditions that prevailed on Twitter before the API changes of 2023, and even more so after. For researchers working on small or under-resourced languages --- where corpora are scarce and institutional resources are limited --- this openness matters. The Slovene community on Bluesky may be small, but it is unusually well-documented as a result.
